\documentclass[10pt,a4paper]{amsart}
\usepackage[margin=19mm]{geometry}
\usepackage{amsmath,amssymb,mathtools}
\usepackage[scr=rsfs,cal=euler]{mathalfa}
\usepackage{xfrac}
\usepackage[unicode,hidelinks,bookmarks=false]{hyperref}
\newcommand{\ie}{i.e.\ }
\newcommand{\cf}{cf.\ }
\newcommand{\resp}{respectively\ }
\newcommand{\Subsection}{\subsection}
\providecommand{\nobreakdash}{\nobreak\hbox{-}}
\setlength{\emergencystretch}{2em}
\multlinegap0pt
\usepackage{float}
\usepackage{amssymb}
\usepackage[normalem]{ulem}
\newtheorem{lemma}{Lemma}
\newtheorem{remark}{Remark}
\newcommand{\bbd}{\mathbb{D}}
\newcommand{\bbn}{\mathbb{N}}
\newcommand{\bbz}{\mathbb{Z}}
\newcommand{\diam}{\text{diam}}
\newcommand{\lc}{\text{lc}}
\newcommand{\mco}{\mathcal{O}}
\newcommand{\ord}{\preceq_{\ast}}
\newcommand{\ov}{\overline}
\newcommand{\scru}{\mathscr{U}}
\newcommand{\ui}{[0,1]}
\title{On R-trees, homotopies, and covering maps}
\author{Jeremy Brazas, Gregory R. Conner, Paul Fabel, Curtis Kent}
\date{Source: arXiv:2401.08883v2 (CC BY 4.0)}
\begin{document}
\maketitle

\noindent\emph{Source and licence.} On R-trees, homotopies, and covering maps. Version arXiv:2401.08883v2 (CC BY 4.0), distributed by arXiv under the Creative Commons Attribution 4.0 International licence, http://creativecommons.org/licenses/by/4.0/, as recorded in the arXiv licence metadata for this exact version and shown on the abstract page https://arxiv.org/abs/2401.08883v2. Full licence notice: \texttt{licenses/arxiv-2401.08883-CC-BY-4.0.txt}.\par\smallskip

\noindent\textit{Editorial scope.} Counted unit: the article's lemma fabelsconstruction (source0.tex lines 520--527). The article's complete original proof of it is delivered with it (source0.tex lines 529--616, one \texttt{proof} environment of 88 lines). No mathematics is written, repaired or re-derived: the statement and the proof are the article's own lines, in the article's own notation, and no retained line is substituted or re-flowed.  Retained uncounted as the source-local context the proof needs: lines 377--379; lines 502--509; lines 511--514.  One declared adaptation: exactly 4 ancillary strings inside the retained spans are omitted (image-inclusion commands and, where the source repeats a label, the repeated label definitions) and no other line differs from the source; the figure environments, with their placement, captions and labels, are kept, and the package records the omitted strings in fidelity receipts bound to the affected bodies.  No retained line contains a citation, so the wrapper carries no bibliography and no bibliography entry.  The retained lines contain the article's own diagram code, which the wrapper typesets with the standard tikz packages; no image file is delivered and the package declares no asset dependency.  Editorial formatting only, in addition: an AMS \texttt{amsart} preamble that restates the article's own declarations and macros used by the retained lines, namely the environments \texttt{lemma}, \texttt{remark} on separate counters, together with the standard packages those lines need.  The retained environments print in this document as Lemma 1, Remark 1 in order of appearance, whereas the article prints the same items with the numbering of its own full document.  The complete native source of the article as distributed in the arXiv e-print for this exact version is bundled unchanged as \texttt{sources/153066/source0.tex}.
\begin{lemma}\label{choosezsetlemma}
Let $U_1$ and $U_2$ be staggered Cantor complements. Then for $i\in\{1,2\}$ there exists a set of connected components $\scru_i$ of $U_i$ such that, $\scru_i$ has the order type of $\bbz$, $\inf(\bigcup\scru_i)=0$, $\sup(\bigcup\scru_i)=1$, and such that $\bigcup\scru_1$ and $\bigcup\scru_2$ are staggered.
\end{lemma}

\begin{remark}\label{linearpathremark}
If $x=y$, then $L_{x,y}$ is constant. If $x\neq y$, then $L_{x,y}$ is a Cantor path that parameterizes the line segment from $x$ to $\frac{x+y}{2}$. Moreover, if $x_1,x_2,y_1,y_2\in \bbd^2$ with midpoints $m_i=\frac{x_i+y_i}{2}$, then since the paths $L_{x_1,y_1}$ and $\ov{L_{y_2,x_2}}$ are both parameterized using $\tau$, their sup-distance is the maximum distance between the endpoints, that is,
\begin{eqnarray*}
\rho(L_{x_1,y_1},\ov{L_{y_2,x_2}}) &\leq& \max\left\{d(x_1,m_2),d(m_1,y_2)\right\}\\
&\leq & \max\left\{\frac{d(x_2,y_2)}{2}+d(x_1,x_2),\frac{d(x_1,y_1)}{2}+d(y_1,y_2)\right\}
\end{eqnarray*}
(See Figure \ref{fig1}). In the case that $x=x_1=x_2$ and $y=y_1=y_2$, we have $\rho(L_{x,y},\ov{L_{y,x}})=\frac{1}{2}d(x,y)$.
\end{remark}

\begin{figure}[H]
\centering 
\caption{\label{fig1}Since $L_{x_1,y_1}$ and $\ov{L_{y_2,x_2}}$ are both parameterized by $\tau$, their metric distance is the maximum length between starting and ending points.}
\end{figure}

\begin{lemma}\label{fabelsconstruction}
For staggered Cantor paths $\alpha,\beta:\ui\to \bbd^2$ such that $\alpha(i)=\beta(i)$, $i\in\{0,1\}$, there exists Cantor paths $\alpha ',\beta ':\ui\to\bbd^2$ such that:
\begin{enumerate}
\item $\alpha\ord \alpha '$ and $\rho(\alpha,\alpha ')\leq \frac{1}{2}\rho(\alpha,\beta)$,
\item $\beta\ord \beta '$ and $\rho(\beta,\beta ')\leq \frac{1}{2}\rho(\alpha,\beta)$,
\item $\rho(\alpha ',\beta ')\leq  \frac{2}{3}\rho(\alpha,\beta)$.
\end{enumerate}
\end{lemma}

\begin{proof}
Let $\delta=\rho(\alpha,\beta)$. By assumption, $\mco(\alpha)$ and $\mco(\beta)$ are staggered Cantor complements. By Lemma \ref{choosezsetlemma}, we may select a set of connected components $\scru_1$ of $\mco(\alpha)$ and $\scru_2$ of $\mco(\beta)$ such that for $i\in\{0,1\}$, $\scru_i$ has the order type of $\bbz$, $\inf(\bigcup\scru_i)=0$, $\sup(\bigcup\scru_i)=1$, and such that $\bigcup\scru_1$ and $\bigcup\scru_2$ are staggered. Index the elements of $\scru_1$ as $\cdots <A_{-2}<A_{-1}<A_0<A_1<A_2<\cdots$ and the elements of $\scru_2$ as $\cdots <C_{-2}<C_{-1}<C_0<C_1<C_2<\cdots$ so that $C_{n}$ meets $A_{n-1}$ and $A_n$. Set $w_n=\alpha(\ov{A_n})$ and $y_n=\beta(\ov{C_n})$

For the moment, fix $n\in\bbz$. Since $\mco(\alpha)$ is a Cantor complement and $\alpha$ is uniformly continuous, we may find a sequence $B_{n,1}<B_{n,2}<\cdots< B_{n,k_n}$ in $\mco(\alpha)$ of length $k_n\geq 2$, where each set $B_{n,j}$ is contained in $[r(A_n), \ell(A_{n+1})]$ and such that if $I$ is a connected component of $[r(A_n), \ell(A_{n+1})]\backslash \bigcup_{j=1}^{k_n}B_{n,j}$, then $\diam(\alpha(\ov{I}))<\frac{\delta}{6}$. Similarly, since $\mco(\beta)$ is a Cantor complement, we may find a sequence $D_{n,1}<D_{n,2}<\cdots< D_{n,m_n}$ in $\mco(\beta)$ of length $m_n\geq 2$ where each $D_{n,j}$ is contained in $[r(C_{n-1}), \ell(C_{n})]$ and such that if $I$ is a connected component of $[r(C_{n-1}), \ell(C_{n})]\backslash \bigcup_{j=1}^{m_n}D_{n,j}$, then $\diam(\beta(\ov{I}))<\frac{\delta}{6}$. Note that $D_{n,j}\subseteq A_n$ and $B_{n,j}\subseteq C_n$ holds whenever these sets are defined (See Figure \ref{fig2}).
\begin{figure}[H]
\centering 
\caption{\label{fig2} The upper segment illustrates the selected sets $A_n$ and $B_{n,j}$ on which $\alpha$ is constant. The lower segment illustrates the selected sets $C_n$ and $D_{n,j}$ on which $\beta$ is constant.}
\end{figure}
Set $x_{n,j}=\alpha(\ov{B_{n,j}})$ and $z_{n,j}=\beta(\ov{D_{n,j}})$ whenever these sets are defined. Additionally, let $p_n$ be the midpoint of $[\ell(A_n),r(C_{n-1})]$, $q_n$ be the midpoint of $[\ell(C_n),r(A_n)]$, $\eta_{n,j}$ be the midpoint of $B_{n,j}$, and $\theta_{n,j}$ be the midpoint of $D_{n,j}$.

To begin our definition of $\alpha'$, we first set $\alpha'$ to agree with $\alpha$ on $\{0,1\}$ and on $[r(A_n),\ell(A_{n+1})]\backslash \bigcup_{j=1}^{k_n}B_{n,j}$ for all $n\in\bbz$.
We complete the definition of $\alpha' $ piecewise by fixing $n\in\bbn$ and defining $\alpha'$ on $\ov{A_n}$ in three cases and $\ov{B_{n,j}}$ for $1\leq j\leq k_n$ in a fourth case. It may be helpful to note that
\[\ov{A_n}=[\ell(A_{n}),p_n]\cup [p_n,r(C_{n-1})]\cup [r(C_{n-1}),\ell(C_n)]\cup [\ell(C_n),q_n]\cup [q_n,r(A_n)]\]
(see Figure \ref{fig3}).

\begin{itemize}
    \item[$\bf (A1)$]\label{alpha'1} On $[\ell(A_{n}),\theta_{n,1}]$, we define $\alpha'$ to be the CIP-loop, which is the linear reparameterization of $L_{w_n,y_{n-1}}$ on $[\ell(A_{n}),p_n]$, the constant path at $\frac{w_n+y_{n-1}}{2}$ on $[p_n,\ell(D_{n,1})]$, and the linear reparameterization of $\ov{L_{w_n,y_{n-1}}}$ on $[\ell(D_{n,1}),\theta_{n,1}]$.
\vspace{.75em}
    \item[$\bf (A2)$]\label{alpha'2} On $[\theta_{n,j},\theta_{n,j+1}]$ (for each $1\leq j\leq m_{n}-1$), we define $\alpha '$ to be the CIP-loop, which is the linear reparameterization of $L_{w_n,z_{n,j}}$ on $[\theta_{n,j},r(D_{n,j})]$, the constant path at $\frac{w_n+z_{n,j}}{2}$ on $[r(D_{n,j}),\ell(D_{n,j+1})]$, and the linear reparameterization of $\ov{L_{w_n,z_{n,j}}}$ on $[\ell(D_{n,j+1}),\theta_{n,j+1}]$.
\vspace{.75em}
    \item[$\bf (A3)$]\label{alpha'3} On $[\theta_{n,m_n},r(A_n)]$, we define $\alpha '$ to be the CIP-loop, which is the linear reparameterization of $L_{w_n, z_{n,m_n}}$ on $[\theta_{n,m_n},r(D_{n,m_n})]$, the constant path at $\frac{w_n+ z_{n,m_n}}{2}$ on $[r(D_{n,m_n}),q_n]$, and the linear reparameterization of $\ov{L_{w_n, z_{n,m_n}}}$ on $[q_n,r(A_n)]$. This completes the definition of $\alpha '$ on $\ov{A_n}$.
\vspace{.75em}
    \item[$\bf (A4)$]\label{alpha'4} Lastly, on $\ov{B_{n,j}}$, we define $\alpha '$ to be the linear reparameterization of the CIP-loop $L_{x_{n,j},y_{n}}\ov{L_{x_{n,j},y_{n}}}$.
\end{itemize}

This completes the definition of $\alpha '$ (compare Figures \ref{fig3} and \ref{figconst}). Note that $\alpha'$ agrees with $\alpha$ everywhere except on a $\bbz$-ordered sequence of elements of $\lc(\alpha)$ on which CIP-loops replace constant loops. Hence, it is clear that $\alpha'$ is a well-defined function.
To begin our definition of $\beta'$, we first set $\beta'$ to agree with $\beta$ on $\{0,1\}$ and on $[r(C_{n-1}),\ell(C_n)]\backslash \bigcup_{j=1}^{m_n}D_{n,j}$ for each $n\in\bbz$. We complete the definition of $\beta' $ piecewise by fixing $n\in\bbn$ and defining $\beta'$ on $\ov{C_n}$ in five cases and $\ov{D_{n,j}}$ for $1\leq j\leq m_n$ in a sixth case.

\begin{itemize}
    \item[$\bf (B1)$]\label{beta'1} On $[\ell(C_n), q_n]$, we define $\beta'$ to be constant at $y_n$ (this happens to agree with the value of $\beta$).
\vspace{.75em}
    \item[$\bf (B2)$]\label{beta'2} On $[ q_n,\eta_{n,1}]$, we define $\beta'$ to be CIP-loops, which is the linear reparameterization of $L_{y_n,w_n}$ on $[ q_n,r(A_n)]$, the constant path at $\frac{y_n+w_n}{2}$ on $[r(A_n),\ell(B_{n,1})]$, and $\ov{L_{y_n,w_n}}$ on $[\ell(B_{n,1}),\eta_{n,1}]$.
\vspace{.75em}
    \item[$\bf (B3)$]\label{beta'3} On $[\eta_{n,j},\eta_{n,j+1}]$ (for each $ 1\leq j\leq k_n-1$), we define $\beta'$ to be CIP-loops, which is the linear reparameterization of $L_{y_n,x_{n,j}}$ on $[\eta_{n,j},r(B_{n,j})]$, the constant path at $\frac{y_n+x_{n,j}}{2}$ on $[r(B_{n,j}),\ell(B_{n,j+1})]$, and the linear reparameterization of $\ov{L_{y_n,x_{n,j}}}$ on $[\ell(B_{n,j+1}),\eta_{n,j+1}]$.
\vspace{.75em}
    \item[$\bf (B4)$]\label{beta'4} On $[\eta_{n,k_n},p_{n+1}]$, we define $\beta'$ to be the CIP-loop, which is the linear reparameterization of $L_{y_n,x_{n,k_n}}$ on $[\eta_{n,k_n},r(B_{n,k_n})]$, the constant path at $\frac{y_n+x_{n,k_n}}{2}$ on $[r(B_{n,k_n}),\ell(A_{n+1})]$, and the linear reparameterization of $\ov{L_{y_n,x_{n,k_n}}}$ on $[\ell(A_{n+1}),p_{n+1}]$.
\vspace{.75em}
    \item[$\bf (B5)$]\label{beta'5} On $[p_{n+1},r(C_n)]$, we define $\beta'$ to be constant at $y_n$ (this happens to agree with the value of $\beta$). This completes the definition of $\beta '$ on $\ov{C_n}$.
\vspace{.75em}
    \item[$\bf (B6)$]\label{beta'6} On $\ov{D_{n,j}}$ (for each $1\leq j\leq m_n$), we define $\beta'$ to be the linear reparameterization of the CIP-loop $L_{z_{n,j},w_{n}} \ov{L_{z_{n,j},w_{n}}}$.
\end{itemize}

This completes the definition of $\beta '$, which is a well-defined function (see Figures \ref{fig3} and \ref{figconst}).
\begin{figure}[H]
\centering 
\caption{\label{fig3} The interlocking pattern that determines the structure of $\alpha '$ and $\beta '$. The triangles represent inserted inverse-pair loops and the trapezoids represent an inverse pair loop but with a constant path included in the middle.  The three subintervals of $\ov{A_n}$ and the five subintervals of $\ov{C_n}$ partitioned by the trapezoids correspond to the piecewise definitions ${\bf(A1)}$-$\bf {(A3)}$ and ${\bf(B1)}$-$\bf {(B5)}$ respectively.}
\end{figure}
\begin{figure}[H]
\centering 
\caption{\label{figconst} A full ``period" of the construction of $\alpha '$ and $\beta '$ starting with $\ell(A_n)$ and ending with $\ell(A_{n+1})$. Starting and ending points are circled and the initial and terminal steps are indicated with arrows. The upper and lower curves represent $\alpha$ and $\beta$ respectively where each subdivided segment has diameter less than $\frac{\delta}{6}$. The numbered paths trace out the trajectory of $\alpha'$ and the box-numbered paths trace out the corresponding trajectory of $\beta '$. When a number is positioned at a point, the path is constant at that point.}
\end{figure}

By construction, $\alpha '|_{(0,1)}$ is continuous. If $U$ is a convex neighborhood of $\alpha(0)=\beta(0)$ in $\bbd^2$, then we may find $N\in\bbz$ such that $\alpha(A_n\cup C_n)\cup \beta(A_n\cup C_n)\subseteq U$ for all $n\leq N$. Let $t=\sup(A_N)$. Since the CIP-loops added to $\alpha$ on $[0,t]$ are contained in line segments with endpoints in $\alpha([0,t])$ and $\beta([0,t])$ respectively, it follows that $\alpha '([0,t])\subseteq U$. Thus $\alpha '$ is continuous at $0$. A symmetric argument shows that $\alpha '$ is continuous at $1$. The construction of $\alpha '$ also ensures that distinct elements of $\lc(\alpha ')$ have disjoint closures. Hence, $\alpha '$ is a Cantor path. Additionally, since $\alpha '$ is constructed from  $\alpha$ only by replacing constant loops with CIP-loops on elements of $\lc(\alpha)$ (one on each $B_{n,j}$ and at least three on each $A_n$), we have $\alpha\ord \alpha '$. Moreover, if $\mu:\ov{I}\to \bbd^2$ is one of the added CIP-loops in the construction of $\alpha '$, then the image of $\mu$ is the line segment connecting $\alpha(s)$ and $\frac{\alpha(s)+\beta(s)}{2}$ for some $s\in \ov{I}$. Thus $\rho(\alpha,\alpha ')\leq \frac{1}{2}\rho(\alpha,\beta)$. Since these arguments apply just as well for $\beta '$, we also conclude that $\beta '$ is a Cantor path satisfying $\beta\ord\beta '$ and $\rho(\beta,\beta ')\leq \frac{1}{2}\rho(\alpha,\beta)$.

To complete the proof, we will show that $d(\alpha '(s),\beta '(s))\leq \frac{2\delta}{3}$. We begin by considering the case $s\in \ov{A_n}$.

\begin{enumerate}
    \item On $[\ell(A_n), p_n]$, $\alpha '$ parameterizes the line from $w_n$ to $\frac{w_n+y_{n-1}}{2}$  (See first part of $\bf (A1)$) and $\beta '$ is a corresponding parameterization of the line from $\frac{x_{n-1, k_{n-1}}+y_{n-1}}{2}$ to $y_{n-1}$  (See last part of $\bf (B4)$). Thus if $s\in [\ell(A_n), p_n]$, Remark \ref{linearpathremark} gives us the first inequality in the following sequence:\\
    \begin{eqnarray*}
     & & d(\alpha '(s),\beta '(s))  \\
    &\leq& \max\left\{\frac{d(x_{n-1, k_{n-1}},y_{n-1})}{2}+d(x_{n-1, k_{n-1}},w_n),\frac{d(w_n,y_{n-1})}{2}\right\}\\
    & =&  \max\left\{\frac{d(\alpha(r(B_{n-1, k_{n-1}})),\beta(r(B_{n-1, k_{n-1}})))}{2}+d(x_{n-1, k_{n-1}},w_n),\frac{d(\alpha(p_n),\beta(p_n))}{2}\right\}\\
    &<& \max\left\{\frac{\delta}{2}+\frac{\delta}{6},\frac{\delta}{2}\right\} \,\leq \, \frac{2\delta}{3}.
    \end{eqnarray*}

    \item On $[p_n,r(C_n)]$, $\alpha '$ is constant at $\frac{w_n+y_{n-1}}{2}$  (See middle part of $\bf (A1)$) and $\beta '$ is constant at $y_{n-1}$  (See $\bf (B5)$). Therefore, if $s\in [p_n,r(C_n)]$, we have \[d(\alpha '(s),\beta '(s))=\frac{d(w_n,y_{n-1})}{2}=\frac{d(\alpha(p_n),\beta(p_n))}{2}\leq\frac{\delta}{2}.\]

    \item On $[r(C_n),\ell(D_{n,1})]$, $\alpha '$ is constant at $\frac{w_n+y_{n-1}}{2}$  (See middle part of $\bf (A1)$) and $\beta '$ agrees with $\beta$. Thus if $s\in [r(C_n),\ell(D_{n,1})]$, then
    \begin{eqnarray*}
    d(\alpha '(s),\beta '(s)) &\leq & d(\alpha '(s),\beta(r(C_{n-1})))+d(\beta(r(C_{n-1})),\beta ' (s))\\
    &= & d\left(\frac{w_n+y_{n-1}}{2},y_{n-1}\right)+d(\beta(r(C_{n-1})),\beta  (s))\\
    &= & \frac{d(\alpha(p_n),\beta(p_n))}{2}+d(\beta(r(C_{n-1})),\beta  (s))\\
    &< & \frac{\delta}{2}+\frac{\delta}{6} \, \leq  \,\frac{2\delta}{3}
    \end{eqnarray*}

    \item On $[\ell(D_{n,1}),\theta_{n,1}]$, $\alpha '$ parameterizes the line segment from $\frac{w_n+y_{n-1}}{2}$ to $w_n$  (See last part of $\bf (A1)$) and $\beta '$ parameterizes the line segment from $z_{n,1}$ to $\frac{z_{n,1}+w_n}{2}$. In this case, if $s\in [\ell(D_{n,1}),\theta_{n,1}]$, then
    \begin{eqnarray*}
    d(\alpha '(s),\beta '(s)) &\leq&  \max\left\{ d\left(w_n,\frac{w_n+y_{n-1}}{2}\right),d\left(z_{n,1},\frac{z_{n,1}+w_n}{2}\right)\right\}\\
    &=& \max\left\{\frac{ d(w_n,y_{n-1})}{2},\frac{d(z_{n,1},w_n)}{2}\right\}\\
    &=& \max\left\{\frac{ d(\alpha(p_n),\beta(p_n))}{2},\frac{d(\alpha(\theta_{n,1}),\beta(\theta_{n,1}))}{2}\right\}\\
    & \leq & \frac{\delta}{2} \,<\, \frac{2\delta}{3}
    \end{eqnarray*}

    \item On $[\theta_{n,1},r(D_{n,1})]$, $\alpha '$ parameterizes the line segment from $w_n$ to $\frac{z_{n,1}+w_n}{2}$  (See first part of $\bf (A2)$) and $\beta$ parameterizes the line segment from $\frac{z_{n,1}+w_n}{2}$ to $z_{n,1}$  (See $\bf (B6)$). Since these paths move along the same line segment, it follows that if $s\in [\theta_{n,1},r(D_{n,1})]$, then \[d(\alpha '(s),\beta '(s))=\frac{d(w_n,z_{n,1})}{2}\leq \frac{d(\alpha(\theta_{n,1}),\beta(\theta_{n,1}))}{2}\leq \frac{\delta}{2}\]
\end{enumerate}

As we proceed through the remainder of the intervals on which $\alpha '$ and $\beta '$ are defined piecewise, every remaining case (all of which are illustrated in Figures \ref{fig3} and \ref{figconst}) may be verified using an argument nearly identical to one of the above five cases. Hence, we omit the remainder of the details. We conclude that $\rho(\alpha',\beta')\leq \frac{2\delta}{3}$.
\end{proof}

\end{document}
